
donc tel que \(\bigoplus_{n\geq n_1} f_*(\underline{F}(n))\) soit un \(\Sym(\underline{E})\)-Module de type fini, donc de présentation finie d'après les hypothèses noethériennes; le résultat est d'ailleurs valable pour tout entier \(n_1\), les \(f_*(\underline{F}(n))\) étant cohérents d'après le théorème de Serre.

La proposition est donc montrée dans le cas où \(S\) est noethérien.

Supposons maintenant qu'on ait un carré cartésien :
\[
\begin{tikzcd}[column sep=large,row sep=large]
X \arrow[r,"g"] \arrow[d,"f"'] & X_0 \arrow[d,"f_0"]\\
S \arrow[r,"g_0"] & S_0,
\end{tikzcd}
\]
avec \(S_0\) noethérien, \(X_0=\mathbf P(\underline{E}_0)\); soit \(\underline{F}_0\) un \(\underline{O}_{X_0}\)-Module de présentation finie, plat sur \(S_0\), et soit \(\underline{F}=g^*(\underline{F}_0)\). D'après ce qui précède la proposition est vraie pour \(\underline{F}_0\). D'après EGA III 7.7.5 et 7.8.5, le théorème de changement de base est vrai avec les hypothèses précédentes; on en déduit aussitôt que le théorème est vrai pour \(\underline{F}\).

Il suffit donc, pour achever la démonstration, de montrer que le cas général ( avec \(S\) affine) peut être obtenu à partir du cas noethérien par changement de base.

Comme \(S\) est affine, \(S\) est limite projective filtrante de schémas affines noethériens (si \(S=\Spec(A)\), on écrit \(A\) comme limite inductive de ses sous-algèbres de type fini sur \(\mathbb Z\)). Soit donc \(S=\varprojlim S_\lambda\). D'après EGA IV 8.5.2 et 8.5.5, il existe un indice \(\lambda\) et un \(\underline{O}_{S_\lambda}\)-Module localement libre de type fini \(\underline{E}_\lambda\) tel que \(\underline{E}=u_\lambda^*(\underline{E}_\lambda)\) (en notant \(u_\lambda\) le morphisme \(S\to S_\lambda\), et \(u_{\lambda\mu}\) les morphismes de transition \(S_\mu\to S_\lambda\)).

Si on pose, pour \(\mu\geq\lambda\), \(\underline{E}_\mu=u_{\lambda\mu}^*(\underline{E}_\lambda)\), et \(X_\mu=\mathbf P(\underline{E}_\mu)\), on a pour tout \(\mu\), \(X=X_\mu\times_{S_\mu}S\), et \(X=\varprojlim X_\mu\). Par suite, il existe un indice \(\lambda'\geq\lambda\)
